\section*{Capítulo \ref{ch:b1:layered-earth} --- La Tierra estratificada}

\begin{solution}{exo:b1:layered-earth:1}
Las fracciones de volumen son los cubos de los cocientes de radios: núcleo $(3480/6371)^3 = 0.1630$, es decir, \qty{16.3}{\percent};
\omterm{def:b1:layered-earth:layers}{núcleo interno} $(1221.5/6371)^3 = 0.00705$, \qty{0.70}{\percent}.
\end{solution}

\begin{solution}{exo:b1:layered-earth:2}
Por encima de la discontinuidad, en la capa superior: $\Phi = 13.7166^2 - \tfrac43\times7.26466^2 = 188.145 - 70.367 =
\qty{117.78}{km^2.s^{-2}}$, $K = \rho\Phi = \qty{5566.45}{kg.m^{-3}}\times
\qty{117.78e6}{m^2.s^{-2}} = \qty{655.6}{GPa}$. Debajo: $\Phi = 8.06482^2 =
\qty{65.04}{km^2.s^{-2}}$, $K = 9903.49\times\num{65.04e6} = \qty{644.1}{GPa}$.
El módulo de compresión apenas cambia (\qty{-1.7}{\percent}). Lo que cambia es el módulo de cizalla,
$\mu = \rho v_S^2 = \qty{293.8}{GPa}$ encima y $0$ debajo (un líquido), y la densidad, que sube un \qty{78}{\percent}:
de ahí la caída de $v_P$ a pesar de un $K$ casi igual.
\end{solution}

\begin{solution}{exo:b1:layered-earth:3}
$\sin i_c = 6.2/8.1 = 0.7654$, $i_c = \ang{49.95}$. Distancia crítica
$2\times30\times\tan\ang{49.95} = 60\times1.1906 = \qty{71.4}{km}$. \omterm{def:b1:layered-earth:head-wave}{Distancia de cruce}
$60\sqrt{14.3/1.9} = 60\times2.743 = \qty{164.6}{km}$.
\end{solution}

\begin{solution}{exo:b1:layered-earth:4}
$g = Gm/r^2 = \num{6.674e-11}\times\num{1.94e24}/(\num{3.48e6})^2 =
\qty{10.69}{m.s^{-2}}$, un \qty{9}{\percent} más que en la superficie. Al bajar al núcleo, el cuadrado del radio cae
hasta $(3480/6371)^2 = 0.298$ de su valor en la superficie, mientras que la masa atractora cae solo hasta
$1.94/5.972 = 0.325$ del total, porque el núcleo contiene un tercio de la masa: $0.325/0.298 = 1.09$.
\end{solution}

\begin{solution}{exo:b1:layered-earth:5}
Cerca de la fuente $v_1 = 60/10.0 = \qty{6.0}{km.s^{-1}}$. Lejos, $v_2 = (300 - 200)/(45.8 - 33.3) = 100/12.5 = \qty{8.0}{km.s^{-1}}$; el tiempo de intersección es $t_i = 33.3 - 200/8.0 = \qty{8.3}{s}$. Entonces $\sin i_c =
0.75$, $\cos i_c = 0.6614$, y $h = t_i v_1/(2\cos i_c) = 8.3\times6.0/1.3229
= \qty{37.6}{km}$. Comprobación: $x_c = 2\times37.6\times\sqrt{14/2} = \qty{199}{km}$, donde ambas ondas llegan en
$199/6.0 = 199/8.0 + 8.3 = \qty{33.2}{s}$.
\end{solution}

\begin{solution}{exo:b1:layered-earth:6}
Para $\rho$ uniforme, $I = \tfrac{8\pi}{3}\rho\int_0^R r^4\,\mathrm dr =
\tfrac{8\pi}{15}\rho R^5$ y $M = \tfrac43\pi\rho R^3$, así que
$I/MR^2 = (8/15)/(4/3) = 0.4$. Dos capas ($r_c = \qty{3.48e6}{m}$,
$R = \qty{6.371e6}{m}$): $M = \tfrac43\pi\bigl(11000\,r_c^3 + 4450\,(R^3 -
r_c^3)\bigr) = \qty{5.977e24}{kg}$, un \qty{0.08}{\percent} por encima de la masa medida; $I = \tfrac{8\pi}{15}\bigl(11000\,r_c^5 + 4450\,(R^5 - r_c^5)\bigr)$ da $I/MR^2 = 0.3457$. La masa es correcta, pero el factor es un \qty{4.5}{\percent} demasiado alto: el \omterm{def:b1:layered-earth:layers}{manto}
real no es uniforme, su densidad crece con la profundidad, lo que lleva la masa hacia dentro.
\end{solution}

\begin{solution}{exo:b1:layered-earth:7}
$\Phi = 12.16^2 - \tfrac43\times6.673^2 = 147.87 - 59.37 =
\qty{88.49}{km^2.s^{-2}}$. Con la profundidad $z$,
$\mathrm d\rho/\mathrm dz = \rho g/\Phi = 4860\times9.9/\num{88.49e6} =
\qty{5.44e-4}{kg.m^{-4}}$, es decir, \qty{0.544}{kg.m^{-3}} por kilómetro, o
\qty{0.54}{g.cm^{-3}} por cada \qty{1000}{km}. El PREM sube $5.357 - 4.844 =
0.513$ en esos \qty{1000}{km}, un \qty{6}{\percent} menos: el gradiente calculado en la parte superior se aplica a todo el
intervalo, mientras que $\Phi$ crece con la profundidad y el gradiente disminuye.
\end{solution}

\begin{solution}{exo:b1:layered-earth:8}
$P \approx \bar\rho g z = 3570\times9.9\times\num{670e3} = \qty{23.7}{GPa}$ (PREM: \qty{23.9}{GPa}). Es la presión, unos
\qty{24}{GPa}, a la que la ringwoodita se descompone en bridgmanita y ferropericlasa: la discontinuidad de \qty{660}{km}.
\end{solution}

\begin{solution}{exo:b1:layered-earth:9}
(a)~Hierro líquido: $10.90\times1.05 = 11.4$ a $10.90\times1.10 =
\qty{12.0}{g.cm^{-3}}$. (b)~$12.764/12.166 = 1.049$: un
salto de \qty{4.9}{\percent} en el PREM, algo más que el \qty{4.5}{\percent} de estudios recientes. (c)~La sola congelación
haría el sólido solo ligeramente más denso que el líquido; el salto mayor significa que el sólido contiene menos elementos
ligeros que el líquido: ellos prefieren el líquido, y la congelación en la \omterm{def:b1:layered-earth:discontinuity}{frontera del núcleo interno} deja la mayor parte atrás.
\end{solution}

\begin{solution}{exo:b1:layered-earth:10}
Sea $x_* = 2h\tan i_c$. Entonces $\sqrt{x_*^2 + 4h^2} = 2h/\cos i_c$ y $t_r(x_*) = 2h/(v_1\cos i_c)$. La \omterm{def:b1:layered-earth:head-wave}{onda de cabeza} da
$t(x_*) = 2h\tan i_c/v_2 + 2h\cos i_c/v_1 = \frac{2h}{v_1}\bigl(\sin^2 i_c/\cos
i_c + \cos i_c\bigr) = 2h/(v_1\cos i_c)$, usando $\sin i_c = v_1/v_2$: lo mismo.
La pendiente de la hipérbola es $t_r'(x) = x/(v_1\sqrt{x^2 + 4h^2})$; en $x_*$ vale
$\sin i_c/v_1 = 1/v_2$, la pendiente de la recta de la \omterm{def:b1:layered-earth:head-wave}{onda de cabeza}. Como
$t_r'' > 0$, la hipérbola es convexa y queda por encima de cada una de sus tangentes: para
$x > x_*$, $t_r > t_{\text{cabeza}}$. En la distancia crítica la reflexión incide sobre la frontera exactamente con $i_c$ y es el
origen de la \omterm{def:b1:layered-earth:head-wave}{onda de cabeza}; más allá, la \omterm{def:b1:layered-earth:head-wave}{onda de cabeza}, que gana por viajar en el medio rápido, llega primero.
\end{solution}

\begin{solution}{exo:b1:layered-earth:11}
$m(r) = \tfrac43\pi\rho r^3$, así que $g = Gm/r^2 = \tfrac43\pi G\rho r$. Entonces
$P(r) = \int_r^R \rho\,\tfrac43\pi G\rho s\,\mathrm ds = \tfrac23\pi G\rho^2(R^2
- r^2)$. En el centro $P = \tfrac23\pi\times\num{6.674e-11}\times5513^2\times
(\num{6.371e6})^2 = \qty{172}{GPa}$, un \qty{47}{\percent} de los
\qty{363.85}{GPa} del PREM: concentrar la masa hacia el centro multiplica por algo más de dos la presión central.
\end{solution}

\begin{solution}{exo:b1:layered-earth:12}
Con $z$ la profundidad, $\mathrm d\rho/\mathrm dz = \rho g/\Phi$; si $g/\Phi$ es una constante $a$,
$\mathrm d(\ln\rho)/\mathrm dz = a$ y $\rho = \rho_0
e^{a(z - z_0)}$. Aquí $a(z - z_0) = 7.9\times\num{2258.5e3}/\num{87e6} =
0.2051$ y $\rho = 9.90\times e^{0.2051} = 9.90\times1.2276 =
\qty{12.15}{g.cm^{-3}}$, un \qty{0.2}{\percent} por debajo de los
$12.17$ del PREM: todo el aumento de densidad a lo largo del \omterm{def:b1:layered-earth:layers}{núcleo externo} es autocompresión.
\end{solution}

\begin{solution}{pb:b1:layered-earth:1}
\textbf{1.} $V = \tfrac43\pi(\num{6.371e6})^3 = \qty{1.0832e21}{m^3}$.

\textbf{2.} $\tfrac43\pi(r_2^3 - r_1^3)$: \omterm{def:b1:layered-earth:layers}{corteza}, \omterm{def:b1:layered-earth:layers}{manto} superior y zona de transición \qty{3.0706e20}{m^3}; \omterm{def:b1:layered-earth:layers}{manto} inferior
\qty{5.9961e20}{m^3}; \omterm{def:b1:layered-earth:layers}{núcleo externo} \qty{1.6890e20}{m^3}; \omterm{def:b1:layered-earth:layers}{núcleo interno} \qty{7.634e18}{m^3}.

\textbf{3.} \qty{28.35}{\percent}, \qty{55.36}{\percent}, \qty{15.59}{\percent} y \qty{0.70}{\percent}: el \omterm{def:b1:layered-earth:layers}{manto} inferior es la
cáscara mayor, más de la mitad del volumen de la Tierra.

\textbf{4.} $15.59 + 0.70 = \qty{16.3}{\percent}$.

\textbf{5.} \omterm{def:b1:layered-earth:layers}{Corteza}, \omterm{def:b1:layered-earth:layers}{manto} superior y zona de transición $\num{3.0706e20}\times3570 =
\qty{1.0962e24}{kg}$; \omterm{def:b1:layered-earth:layers}{manto} inferior
\qty{2.9381e24}{kg}; \omterm{def:b1:layered-earth:layers}{núcleo externo} \qty{1.8410e24}{kg}; \omterm{def:b1:layered-earth:layers}{núcleo interno} \qty{9.841e22}{kg}.

\textbf{6.} \qty{5.9737e24}{kg}, un \qty{0.03}{\percent} por encima de la masa medida \qty{5.972e24}{kg}.

\textbf{7.} $\num{5.9737e24}/\num{1.0832e21} = \qty{5515}{kg.m^{-3}}$.

\textbf{8.} $\num{9.841e22}/\num{5.9737e24} = \qty{1.65}{\percent}$.

\textbf{9.} $(1.0962 + 2.9381)/5.9737 = \qty{67.5}{\percent}$.

\textbf{10.} El núcleo es el \qty{16.3}{\percent} del volumen y el \qty{32.5}{\percent} de la masa, el doble: su densidad media,
\qty{10986}{kg.m^{-3}}, es unas $2.5$ veces la del \omterm{def:b1:layered-earth:layers}{manto}.

\textbf{11.} Una cáscara de densidad uniforme da
$I = \tfrac{8\pi}{3}\rho\int_{r_1}^{r_2} r^4\,\mathrm dr = \tfrac{8\pi}{15}\rho
(r_2^5 - r_1^5)$.

\textbf{12.} \omterm{def:b1:layered-earth:layers}{Corteza}, \omterm{def:b1:layered-earth:layers}{manto} superior y zona de transición \qty{2.676e37}{kg.m^2}; \omterm{def:b1:layered-earth:layers}{manto} inferior \qty{4.525e37}{kg.m^2};
\omterm{def:b1:layered-earth:layers}{núcleo externo} \qty{9.27e36}{kg.m^2}; \omterm{def:b1:layered-earth:layers}{núcleo interno} \qty{5.87e34}{kg.m^2}.

\textbf{13.} El \omterm{def:b1:layered-earth:layers}{manto} inferior: tiene la mitad de la masa, a radios grandes. El \omterm{def:b1:layered-earth:layers}{núcleo interno} contiene el \qty{1.6}{\percent} de la masa,
y toda ella cerca del eje: su parte es del \qty{0.07}{\percent}.

\textbf{14.} $I = \qty{8.1345e37}{kg.m^2}$ e $I/MR^2 = \num{8.1345e37}/
(\num{5.9737e24}\times(\num{6.371e6})^2) = 0.3355$.

\textbf{15.} Un \qty{1.4}{\percent} por encima de $0.3308$, y bastante por debajo de $0.4$: el modelo tiene casi la
concentración correcta de masa hacia el centro.

\textbf{16.} Sí. Dentro de cada cáscara la densidad real es mayor abajo que arriba; una cáscara uniforme de la misma masa pone
demasiada masa cerca de su radio exterior, donde cuenta con $r^2$, y por tanto sobrestima $I$.

\textbf{17.} $g = GM/R^2 = \num{6.674e-11}\times\num{5.9737e24}/
(\num{6.371e6})^2 = \qty{9.822}{m.s^{-2}}$.

\textbf{18.} La masa del núcleo es $\num{1.8410e24} + \num{9.841e22} =
\qty{1.9394e24}{kg}$, y
$g = G m/r_c^2 = \qty{10.69}{m.s^{-2}}$.

\textbf{19.} $(\num{1.0962e24} + \num{2.9381e24})/(\num{3.0706e20} +
\num{5.9961e20}) = \num{4.0343e24}/\num{9.0667e20} = \qty{4450}{kg.m^{-3}}$.

\textbf{20.} $P \approx 4450\times10\times\num{2.891e6} = \qty{128.6}{GPa}$, un \qty{5}{\percent} por debajo de los
\qty{135.75}{GPa} del PREM.

\textbf{21.} $\rho_c = \num{1.9394e24}/\num{1.7653e20} = \qty{10986}{kg.m^{-3}}$, y $\Delta P = \tfrac23\pi\times\num{6.674e-11}\times10986^2\times
(\num{3.48e6})^2 = \qty{204.3}{GPa}$.

\textbf{22.} $128.6 + 204.3 = \qty{332.9}{GPa}$, un \qty{8.5}{\percent} por debajo de los \qty{363.85}{GPa} del PREM. Ambas
estimaciones ignoran el aumento de densidad con la profundidad dentro de cada región, que sitúa más masa en profundidad, mantiene la
gravedad más alta al bajar y aumenta $\rho g$, el gradiente de presión.

\textbf{23.} La masa total (pregunta 6), el \omterm{def:b1:comparative-planetology:moi}{factor de momento de inercia} (pregunta 15) y la gravedad en la superficie (pregunta 17,
frente a la media medida de \qty{9.82}{m.s^{-2}}) se comprueban con medidas de la Tierra; la gravedad en el núcleo y las presiones
solo pueden compararse con el PREM, que es en sí mismo un modelo.

\textbf{24.} El núcleo contiene el \textbf{\qty[text-series-to-math=true]{32.5}{\percent}} de la masa de la Tierra.
\end{solution}
